%-----------------------------------------------------------------------------%
\chapter{\babSatu}
\label{cha:pendahuluan}
%-----------------------------------------------------------------------------%
Bab ini memaparkan alasan ilmiah dan teknis penelitian DroneVLM. Pembahasan dimulai dari kebutuhan navigasi UAV yang memahami tujuan semantik, kemudian menguraikan masalah integrasi VLM yang heterogen pada sistem navigasi \textit{closed-loop}. Selanjutnya, bab ini merumuskan masalah dan tujuan penelitian, menjelaskan manfaat serta batasan penelitian, dan menempatkan setiap bab sebagai bagian dari alur pembuktian penelitian.

%-----------------------------------------------------------------------------%
\section{Latar Belakang}
\label{sec:latar_belakang}
%-----------------------------------------------------------------------------%
Wahana udara nirawak atau \textit{unmanned aerial vehicle} (UAV) menjalankan misi pada lingkungan yang informasi relevannya dapat berubah selama penerbangan. Pada kondisi tersebut, perencanaan lintasan tidak hanya menentukan rute dari satu koordinat ke koordinat lain, tetapi juga memperbarui keputusan berdasarkan pengamatan lingkungan dan keadaan kendaraan. Tinjauan Vivaldini dkk. menempatkan pengambilan keputusan dari informasi yang diperoleh selama penerbangan sebagai bagian penting dari perencanaan lintasan UAV dalam lingkungan yang tidak pasti \citep{vivaldini2025decision}. Oleh karena itu, navigasi UAV untuk misi semantik memerlukan mekanisme yang dapat menghubungkan pengamatan terkini dengan keputusan gerak berikutnya. % [SRC-B1-01]

Navigasi berbasis tujuan-objek memperluas masalah tersebut karena sasaran tidak selalu tersedia sebagai koordinat yang telah diketahui. UAV perlu mencari objek, mengenali kandidat yang sesuai dengan instruksi, dan memperkirakan lokasi sasaran sebelum membentuk subtujuan geometris yang aman. Ayala dkk. menjelaskan bahwa pencarian dan identifikasi objek merupakan bagian sentral dari navigasi tujuan-objek pada UAV dan berkaitan dengan perencanaan, pelokalan, serta kendali \citep{ayala2024uav}. Dengan demikian, sistem tidak cukup hanya mengoptimalkan lintasan; sistem juga harus memahami apa yang dicari dan bukti visual apa yang mendukung keputusan menuju sasaran tersebut. % [SRC-B1-02]

Instruksi bahasa alami menambah kebutuhan untuk menambatkan entitas bahasa ke lingkungan visual. Agen harus menghubungkan objek, atribut, dan relasi seperti kiri, kanan, depan, atau di antara dengan pengamatan kamera; agen juga harus mengetahui subinstruksi yang telah diselesaikan dan subinstruksi yang perlu dilakukan berikutnya. Penelitian navigasi visi-bahasa menunjukkan bahwa penambatan landmark serta relasi spasial ke modalitas visual penting bagi pelaksanaan instruksi, sedangkan pemantauan kemajuan membantu agen memilih tindakan berikutnya secara konsisten \citep{zhang2022explicit,ma2019selfmonitoring}. Kebutuhan tersebut menjadi lebih menantang pada UAV karena perubahan yaw, ketinggian, dan sudut pandang dapat mengubah bukti visual yang tersedia. % [SRC-B1-03]

Model visi-bahasa atau \textit{vision-language model} (VLM) berpotensi menyediakan penalaran tingkat tinggi yang diperlukan untuk menginterpretasikan pengamatan dan instruksi secara bersama-sama. CLIP menunjukkan bahwa representasi gambar--teks dapat dipelajari dari supervisi bahasa alami dan ditransfer ke konsep visual di luar kelas objek tetap \citep{radford2021clip}. Sementara itu, penelitian pembangkitan keterangan gambar menunjukkan keterhubungan visi komputer dan pemrosesan bahasa alami dalam menghasilkan deskripsi berbasis gambar \citep{vinyals2015show}. Meskipun demikian, keluaran VLM bersifat probabilistik dan dapat berbentuk bahasa bebas, sehingga keluaran tersebut belum dapat diperlakukan sebagai perintah penerbangan tanpa normalisasi dan pemeriksaan lebih lanjut. % [SRC-B1-04]

Masalah integrasi muncul karena setiap \textit{backend} VLM dapat memiliki kebutuhan resolusi gambar, prapemrosesan, templat \textit{prompt}, struktur konteks, dan format respons yang berbeda. Satu model dapat mengembalikan teks bebas, sedangkan model lain dapat mengembalikan pilihan diskret atau objek JSON. Apabila simpul simulator, pengurai respons, dan logika penerbangan ditulis khusus untuk satu \textit{backend}, penggantian model akan mengubah jalur kendali sekaligus. Perbandingan yang dihasilkan menjadi sulit ditafsirkan karena perbedaan kinerja dapat berasal dari perubahan pengurai atau pengendali, bukan dari kemampuan model. Penelitian ini karena itu memerlukan batas antarmuka yang memisahkan inferensi spesifik-model dari simulasi dan kendali yang tetap.

Pemisahan tersebut penting dalam navigasi \textit{closed-loop}, yaitu ketika keputusan saat ini mengubah pengamatan pada langkah berikutnya. Respons VLM yang tidak sesuai skema, ambigu, memakai satuan yang tidak didukung, berasal dari data sensor yang kedaluwarsa, atau mengarahkan UAV ke rintangan tidak boleh langsung dieksekusi. DroneVLM menempatkan VLM sebagai perencana semantik tingkat tinggi, bukan pengendali aktuator. Kontrak pengamatan, pengurai aksi, gerbang keselamatan, pengendali tingkat rendah, mesin keadaan misi, dan pencatat eksperimen dipertahankan sebagai komponen yang terpisah dari \textit{backend}. Rancangan ini membatasi dampak satu respons yang buruk dan memungkinkan sumber kegagalan dicatat secara terpisah.

CoSyS-AirSim digunakan sebagai lingkungan simulasi karena konfigurasi penelitian dapat menyediakan kendaraan multirotor, citra RGB, kedalaman, segmentasi instans, dan lapisan anotasi. Lingkungan tersebut dihubungkan dengan ROS~2 untuk pertukaran citra kamera, odometri, IMU, GPS, data jarak, transformasi objek, perintah kecepatan, dan tujuan posisi lokal. ROS~2 menyediakan abstraksi simpul, topik, layanan, aksi, parameter, serta transformasi yang mendukung integrasi komponen robotika \citep{macenski2022ros2}. Dalam penelitian ini, data anotasi simulator digunakan untuk membangun kebenaran dasar dan mengevaluasi misi, bukan sebagai masukan utama bagi VLM. % [SRC-B1-05]

Repositori dasar penelitian telah menyediakan CoSyS-AirSim, paket jembatan ROS~2, konfigurasi multirotor \texttt{SimpleFlight}, dan pengendali posisi PD sederhana. Namun, repositori tersebut belum menyediakan adaptor VLM, pengelola episode, evaluator arena, maupun \textit{backend} VLM yang siap dibandingkan. Penelitian ini mengisi batas tersebut melalui DroneVLM dan DroneVLM-Arena. DroneVLM menormalkan pengamatan dan respons model, sedangkan DroneVLM-Arena mendefinisikan episode, aturan terminal, anotasi, log, serta evaluator. Dengan pembagian ini, perubahan model tidak memerlukan perubahan pada pengendali atau prosedur evaluasi.

Satu demonstrasi UAV yang mencapai sasaran belum cukup untuk mendukung klaim tentang kemampuan navigasi semantik. Evaluasi harus menggunakan episode dengan pose awal, instruksi, sasaran, batas waktu, anggaran keputusan, dan batasan keselamatan yang ditetapkan terlebih dahulu. Selain tingkat keberhasilan misi, evaluasi perlu mencatat tabrakan, aksi tidak valid, intervensi keselamatan, panjang lintasan, latensi inferensi, dan latensi keputusan ujung-ke-ujung. Benchmark VLM untuk UAV juga menekankan perlunya evaluasi kecerdasan spasial yang mempertimbangkan perspektif serta karakteristik navigasi UAV, bukan hanya jawaban statis dari perspektif manusia \citep{zhang2025skyready}. % [SRC-B1-06]

Berdasarkan kebutuhan tersebut, penelitian ini mengusulkan \textbf{DroneVLM: Adaptor VLM Agnostik-Model untuk Navigasi UAV Semantik di CoSyS-AirSim}. Kontribusi penelitian terdiri atas kontrak pengamatan dan aksi bersama, pemisahan adaptor dari jalur kendali, gerbang keselamatan operasional untuk simulasi, serta arena evaluasi yang dapat direproduksi. Penelitian ini tidak mengusulkan VLM baru dan tidak membuat klaim keselamatan penerbangan di dunia nyata. Hasil yang dilaporkan nantinya dibatasi pada model, adegan, konfigurasi sensor, kondisi lingkungan, dan prosedur evaluasi yang benar-benar diuji dalam DroneVLM-Arena.

%-----------------------------------------------------------------------------%
\section{Rumusan Masalah}
\label{sec:rumusan_masalah}
%-----------------------------------------------------------------------------%
Berdasarkan latar belakang tersebut, rumusan masalah penelitian adalah sebagai berikut.
\begin{enumerate}
    \item Bagaimana merancang dan mengintegrasikan VLM yang heterogen ke dalam satu kerangka navigasi UAV semantik tanpa modifikasi spesifik-model pada tumpukan kendali hilir?
    \item Bagaimana mengevaluasi dan membandingkan VLM pada tugas penambatan visual, penalaran spasial, pencarian objek, navigasi semantik, dan pelaksanaan misi multilangkah dalam lingkungan simulasi yang terstandar?
    \item Sejauh mana adaptasi \textit{supervised fine-tuning} hemat-parameter menggunakan LoRA dan \textit{post-training} berbasis penguatan menggunakan GRPO dapat meningkatkan kinerja navigasi \textit{closed-loop} VLM tujuan umum?
    \item Bagaimana kompromi antara ukuran model, \textit{backend} inferensi, latensi, penggunaan sumber daya komputasi, efisiensi lintasan, keselamatan, dan keberhasilan misi?
    \item Seberapa tangguh perencana UAV berbasis VLM terhadap variasi lingkungan, perumusan \textit{prompt}, modalitas pengamatan, kondisi visual, dan kompleksitas misi; serta bagaimana generalisasinya pada skenario yang belum dilihat?
\end{enumerate}

%-----------------------------------------------------------------------------%
\section{Tujuan Penelitian}
\label{sec:tujuan_penelitian}
%-----------------------------------------------------------------------------%
Selaras dengan rumusan masalah, penelitian ini bertujuan untuk:
\begin{enumerate}
    \item merancang dan mengimplementasikan DroneVLM sebagai adaptor VLM agnostik-model pada CoSyS-AirSim dan ROS~2;
    \item mendefinisikan kontrak pengamatan dan aksi yang memisahkan persiapan serta inferensi spesifik-model dari simulasi, kendali, dan evaluasi;
    \item mengimplementasikan pengurai, penormal aksi, dan gerbang keselamatan untuk memvalidasi aksi tingkat tinggi sebelum dieksekusi;
    \item merancang DroneVLM-Arena yang memuat episode misi berversi, kebenaran dasar semantik, aturan terminal, pencatatan, dan evaluator;
    \item mengevaluasi dan membandingkan VLM terpilih pada penambatan visual, penalaran spasial, pencarian objek, navigasi semantik, dan instruksi multilangkah;
    \item menganalisis keberhasilan misi, keselamatan, efisiensi lintasan, aksi tidak valid, penggunaan sumber daya, dan latensi; serta
    \item mengevaluasi pengaruh adaptasi LoRA dan GRPO terhadap kinerja navigasi \textit{closed-loop} apabila data pelatihan dan sumber daya eksperimen tersedia.
\end{enumerate}

%-----------------------------------------------------------------------------%
\section{Manfaat Penelitian}
\label{sec:manfaat_penelitian}
%-----------------------------------------------------------------------------%
Penelitian ini diharapkan memberikan manfaat sebagai berikut.
\begin{enumerate}
    \item \textbf{Bagi pengembang UAV dan robotika.} DroneVLM menyediakan batas yang jelas antara \textit{backend} VLM, simulator, validasi keselamatan, dan pengendali penerbangan. Pengembang dapat menguji \textit{backend} baru melalui adaptor tanpa menulis ulang jalur ROS~2, pengendali tingkat rendah, atau evaluator.
    \item \textbf{Bagi peneliti \textit{embodied AI}.} DroneVLM-Arena menyediakan rancangan evaluasi \textit{closed-loop} dengan episode dan prosedur pencatatan yang tetap. Rancangan tersebut membantu membedakan kegagalan inferensi, pengurai, keselamatan, pengendali, dan infrastruktur dari sekadar hasil akhir misi.
    \item \textbf{Bagi pengguna akademik.} Penelitian ini menyediakan contoh integrasi VLM--ROS~2 pada UAV simulasi dan dasar untuk penelitian lanjutan tentang perencanaan tangguh, keselamatan operasional, serta \textit{sim-to-real transfer}. Hasil evaluasi dapat digunakan untuk menentukan batas kemampuan sistem pada kondisi yang diuji.
    \item \textbf{Bagi perancang sistem.} Analisis metrik kualitas, keselamatan, efisiensi, sumber daya, dan latensi dapat menjadi dasar pemilihan \textit{backend} sesuai kebutuhan misi. Pemilihan tersebut tidak hanya mempertimbangkan keberhasilan misi, tetapi juga biaya dan risiko operasional pada jalur keputusan.
\end{enumerate}

%-----------------------------------------------------------------------------%
\section{Ruang Lingkup Penelitian}
\label{sec:ruang_lingkup}
%-----------------------------------------------------------------------------%
Ruang lingkup penelitian dibatasi sebagai berikut.
\begin{enumerate}
    \item Evaluasi dilakukan pada CoSyS-AirSim Blocks. Penelitian tidak mencakup penerbangan dunia nyata, pengujian \textit{hardware-in-the-loop}, maupun klaim keselamatan penerbangan nyata.
    \item Kendaraan yang digunakan adalah satu UAV multirotor dalam mode \texttt{Multirotor} dengan \texttt{SimpleFlight}. Pengendali tingkat rendah dibuat tetap pada seluruh perbandingan \textit{backend}.
    \item VLM berperan sebagai perencana semantik tingkat tinggi dan tidak dapat menerbitkan perintah motor atau gaya dorong secara langsung.
    \item Pengamatan utama terdiri atas citra RGB depan, kedalaman terdaftar, keadaan atau odometri UAV, ringkasan kemajuan misi, dan instruksi bahasa alami. LiDAR, fusi banyak kamera, SLAM penuh, dan memori jangka panjang berada di luar konfigurasi utama kecuali sebagai ablasi.
    \item \textit{Backend} utama yang dikunci adalah LLaVA-v1.5-7B, Qwen2-VL-2B-Instruct, dan InternVL2-2B. Revisi model, checksum SHA-256 bobot, kuantisasi, \textit{runtime}, prosesor, \textit{prompt}, serta parameter dekode dicatat sebelum pengujian.
    \item Setiap \textit{backend} menggunakan skema aksi, pengendali, batas gerak, anggaran waktu dan keputusan, geofence, gerbang keselamatan, kebijakan coba-ulang, serta evaluator yang sama. Perbedaan kemampuan bawaan dilaporkan sebagai hasil, tidak dinormalisasi atau disembunyikan.
    \item Benchmark dibatasi pada penambatan visual, penalaran spasial, pencarian objek, navigasi semantik, dan instruksi multilangkah. Manipulasi, koordinasi banyak UAV, pendaratan presisi, dan interaksi fisik tidak dibahas.
    \item Label simulator dan transformasi objek hanya digunakan untuk anotasi, baseline oracle, dan evaluasi; keduanya tidak diberikan kepada VLM pada kondisi utama.
    \item Kesimpulan dibatasi pada adegan, konfigurasi sensor, distribusi bahasa, model, serta kondisi lingkungan yang ditetapkan dan diuji dalam DroneVLM-Arena.
\end{enumerate}

%-----------------------------------------------------------------------------%
\section{Sistematika Penulisan}
\label{sec:sistematika}
%-----------------------------------------------------------------------------%
Laporan ini disusun dalam lima bab yang saling berkaitan.
\begin{enumerate}
    \item \textbf{Bab I Pendahuluan} memaparkan konteks masalah, rumusan masalah, tujuan, manfaat, ruang lingkup, dan alur penulisan penelitian.
    \item \textbf{Bab II Tinjauan Pustaka} menjelaskan teori yang mendasari DroneVLM, penelitian terkait, celah penelitian, dan posisi kontribusi yang diusulkan.
    \item \textbf{Bab III Metodologi dan Perancangan Sistem} menjabarkan desain penelitian, rancangan arsitektur, kontrak antarkomponen, episode evaluasi, metrik, dan protokol pencatatan.
    \item \textbf{Bab IV Pengujian dan Evaluasi} menyajikan konfigurasi eksperimen, hasil, validasi data, dan analisis keberhasilan misi, keselamatan, efisiensi, serta latensi berdasarkan log evaluator.
    \item \textbf{Bab V Penutup} merangkum jawaban atas rumusan masalah berdasarkan bukti Bab IV dan menyajikan saran pengembangan penelitian berikutnya.
\end{enumerate}
